% Thesis Abstract -----------------------------------------------------


%\begin{abstractslong}    %uncommenting this line, gives a different abstract heading
\begin{abstracts}        %this creates the heading for the abstract page

Scientific workflows underpin modern data-driven research, yet composing them
remains a task reserved for those with programming expertise. Domain scientists
like clinicians, biologists, policy analysts, must all either learn a
workflow-description language or rely on software engineers to express their
analytical intent as an executable pipeline. With the recent emergence of large 
language models (LLMs), an opportunity to bridge this gap by translating natural-language
descriptions of analytical intent into executable workflows arises.
At the same time, relying on LLMs to automate this translation 
raises two practical concerns that are rarely addressed together: 
\emph{privacy}, since scientific data is often sensitive and cannot be sent to a commercial cloud API, 
and \emph{reproducibility}, since a system that generates a different workflow each
time the same question is asked undermines the very goal of scientific
computing.

This thesis explores whether these constraints can be satisfied simultaneously.
It proposes a pipeline built entirely on locally deployable, open-weight models
that translates natural-language descriptions of analytical intent into
executable workflows for the Brane distributed computing framework. By keeping
inference on-premises and caching verified translations by semantic
equivalence, the system attempts to close the gap between the ease of
natural-language expression and the precision that scientific workflow
execution demands, without compromising on data governance or result
consistency. The findings offer insight into where the real bottlenecks lie
when bringing LLM-based automation into resource-constrained, privacy-sensitive
scientific environments.

The pipeline combines retrieval-augmented generation, structured intent
decomposition, and a semantic cache, evaluated across three open-weight model
scales on a BraneScript test set. The results show that prompt and retrieval
engineering outweigh model scale as the primary driver of translation quality,
with the largest model reaching 93\% execution rate after targeted grounding
improvements. Intent decomposition provides consistent gains at every scale.
Supervised fine-tuning revealed a prompt-coupling fragility that must be
managed explicitly in any maintained deployment. The semantic cache achieves
near-complete hit rates on paraphrased inputs while exposing a fundamental
limit of embedding-based similarity for parameter-sensitive scientific queries
--- a tension that points to open problems at the intersection of NLP and
scientific reproducibility.

\end{abstracts}
%\end{abstractlongs}


% ----------------------------------------------------------------------
